% ID: SEL-D0042-OB-E01
% ESTADO: cotejada. Segunda pasada visual y borrador compilado revisados.
% Fuente: Selectividad/sel_2014_matematicas_aplicadas/EXAMEN -ANDAL-1.pdf
% SHA-256: 375424f4e0d3c0409619ebaae4ae4d20c441a2585efe3dc74eb98d765114311e
% Página: 2; opción B; ejercicio original 1.

Un nutricionista receta a una de sus pacientes una dieta semanal especial basada en lácteos y pescado. Cada kg de lácteos cuesta 6 \texteuro{} y proporciona 3 unidades de proteínas y 1 de calorías; cada kg de pescado cuesta 12 \texteuro{}, aportando 1 unidad de proteínas y 2 de calorías.

La dieta le exige no tomar más de 4 kg, conjuntamente, de lácteos y pescado, y un aporte mínimo de 4 unidades de proteínas y 3 de calorías.
\begin{enumerate}[label=\alph*)]
\item \textbf{(1 punto)} Plantee el problema para obtener la combinación de ambos alimentos que tenga el coste mínimo.
\item \textbf{(1.5 puntos)} Dibuje la región factible y determine la solución óptima del problema.
\end{enumerate}
